\begin{enumerate}
\item \begin{enumerate}
\item On a :
\[
\vect{BC}\begin{pmatrix}3-1\\3-2\\1-3\end{pmatrix}\quad \text{soit}\quad \vect{BC}\begin{pmatrix}2\\1\\-2\end{pmatrix}
\]
Et :
\[
\vect{BE}\begin{pmatrix}2-1\\-2-2\\2-3\end{pmatrix}\quad \text{soit}\quad \vect{BE}\begin{pmatrix}1\\-4\\-1\end{pmatrix}
\]
Ces coordonnées ne sont pas proportionnelles, donc les vecteurs ne sont pas colinéaires, et donc les trois points B, C et E ne sont pas alignés.

Ils définissent donc un plan $(BCE)$.

\item On a :
\[
\vect{AF}\begin{pmatrix}3-0\\0-0\\4-1\end{pmatrix}\quad \text{soit}\quad \vect{AF}\begin{pmatrix}3\\0\\3\end{pmatrix}
\]
Comme le repère est orthonormé, on va utiliser la formule analytique pour calculer les produits scalaires :
\begin{itemize}
\item $\vect{AF}\cdot \vect{BC}=3\times 2+0\times 1+3\times(-2)=6+0-6=0$.

$\vect{AF}$ et $\vect{BC}$ sont orthogonaux;
\item $\vect{AF}\cdot \vect{BE}=3\times 1+0\times(-4)+3\times(-1)=3+0-3=0$.

$\vect{AF}$ et $\vect{BE}$ sont orthogonaux.
\end{itemize}
$\vect{AF}$ est donc orthogonal à deux vecteurs non colinéaires qui forment une base du plan $(BCE)$, donc il est normal au plan $(BCE)$.

\item Par propriété, on en déduit qu'une équation cartésienne du plan $(BCE)$ est de la forme $3\times x+0\times y+3\times z+d=0$.

B étant un point du plan $(BCE)$, ses coordonnées doivent vérifier l'équation du plan.
\[
B\in (BCE)\Longleftrightarrow 3\times 1+3\times 3+d=0
\Longleftrightarrow 12+d=0
\Longleftrightarrow d=-12.
\]
Une équation cartésienne du plan est donc $3x+3z-12=0$.

En divisant cette équation par 3 (qui est non nul), on obtient une équation équivalente, donc une équation du plan $(BCE)$ est bien : $x+z-4=0$.
\end{enumerate}

\item \begin{enumerate}
\item Voyons si les coordonnées de G vérifient l'équation du plan :
\[
x_G+z_G-4=5+1-4=2\neq 0.
\]
Les coordonnées de G ne vérifient pas l'équation du plan $(BCE)$ donc le point G n'appartient pas au plan $(BCE)$.

\item Comme $\vect{BE}$ et $\vect{BC}$ ne sont pas colinéaire (prouvé à la question 1. a.), on en déduit que les vecteurs $\vect{BE}$, $\vect{BC}$ et $\vect{AG}$ sont coplanaires si et seulement si on peut exprimer $\vect{AG}$ comme combinaison linéaire des deux autres vecteurs.

Autrement dit, si et seulement s'il existe deux réels $a$ et $b$ tels que : $\vect{AG}=a\vect{BC}+b\vect{BE}$.

On a :
\[
\vect{AG}\begin{pmatrix}5\\1\\1\end{pmatrix}.
\]
Donc :
\[
\begin{aligned}
\vect{AG}=a\vect{BC}+b\vect{BE}
&\Longleftrightarrow
\begin{cases}
2a+b=5\\
a-4b=1\\
2a-b=1
\end{cases}\\
&\Longleftrightarrow
\begin{cases}
b=5-2a\\
a-4(5-2a)=1\\
2a-(5-2a)=1
\end{cases}\\
&\Longleftrightarrow
\begin{cases}
b=5-2a\\
9a-20=1\\
4a-5=1
\end{cases}\\
&\Longleftrightarrow
\begin{cases}
b=5-2a\\
9a=21\\
4a=6
\end{cases}\\
&\Longleftrightarrow
\begin{cases}
b=5-2a\\
a=\dfrac73\\
a=\dfrac32.
\end{cases}
\end{aligned}
\]
Les équations du système ne sont pas compatibles.

Donc le système n'a pas de solution.

Donc les vecteurs $\vect{BE}$, $\vect{BC}$ et $\vect{AG}$ ne sont pas coplanaires.

\item La droite $(AG)$ est dirigée par $\vect{AG}$, qui n'est pas un vecteur du plan $(BCE)$, donc la droite $(AG)$ n'est pas parallèle au plan (ni strictement parallèle, ni incluse dans le plan).

Une droite et un plan n'ont que trois positions relatives possibles, et on vient de justifier que deux ne sont pas possibles.

Par élimination, on peut conclure que la droite $(AG)$ et le plan $(BCE)$ sont sécants.
\end{enumerate}

\item \begin{enumerate}
\item Par propriété, une représentation paramétrique de $(AG)$ est :
\[
\begin{cases}
x=x_A+x_{\vect{AG}}\times t\\
y=y_A+y_{\vect{AG}}\times t\\
z=z_A+z_{\vect{AG}}\times t
\end{cases}\quad t\in\R.
\]
En remplaçant par le coordonnées connues, cela donne :
\[
\begin{cases}
x=0+5\times t\\
y=0+1\times t\\
z=1+1\times t
\end{cases}\quad t\in\R.
\]
On a donc bien une représentation paramétrique de $(AG)$ qui est :
\[
\begin{cases}
x=5t\\
y=t\\
z=1+t
\end{cases}\quad t\in\R.
\]

\item Par définition, le point P est un point de $(AG)$.

Ses coordonnées sont donc de la forme $(5t~;~t~;~1+t)$ avec un $t$ réel.

Comme c'est aussi un point du plan $(BCE)$, ces coordonnées doivent vérifier l'équation du plan :
\[
x_P+z_P-4=0\Longleftrightarrow 5t+(1+t)-4=0
\Longleftrightarrow 6t=3
\Longleftrightarrow t=0,5.
\]
P a donc pour coordonnées : $(5\times 0,5~;~0,5~;~1+0,5)$ soit $(2,5~;~0,5~;~1,5)$.

\item Le milieu du segment $[EC]$ a pour coordonnées :
\[
\left(\dfrac{x_E+x_C}{2}~;~\dfrac{y_E+y_C}{2}~;~\dfrac{z_E+z_C}{2}\right).
\]
En remplaçant par les valeurs connues :
\[
\left(\dfrac{2+3}{2}~;~\dfrac{-2+3}{2}~;~\dfrac{2+1}{2}\right).
\]
En effectuant les calculs, on retrouve $(2,5~;~0,5~;~1,5)$ c'est-à-dire les coordonnées de P.

P est donc bien le milieu de $[EC]$.
\end{enumerate}

\item Les plans $(BCE)$ et $(ACG)$ sont sécants, en effet :
\begin{itemize}
\item ils ont le point C en commun, de façon évidente : on en déduit que les plans ne sont pas strictement parallèles;
\item ils ne sont pas confondus, car le point G appartient à $(ACG)$, mais pas à $(BCE)$ (prouvé à la question 2. a.).
\end{itemize}
Deux plans sécants ont une intersection qui est une droite, que l'on va nommer $d$.
\begin{itemize}
\item P étant par définition sur la droite $(AG)$, il est un point de $(ACG)$;
\item P est aussi par définition un point du plan $(BCE)$.
\end{itemize}
P est donc un autre point commun aux plans $(ACG)$ et $(BCE)$, donc la droite $(CP)$ est incluse dans l'intersection des plans $(ACG)$ et $(BCE)$.

Comme cette intersection est réduite à une droite, on en déduit que :

l'intersection des plans $(BCE)$ et $(ACG)$ est la droite $(CP)$.
\end{enumerate}
